\documentclass[11pt]{article}

% ===================== PACKAGES =====================
\usepackage[a4paper,margin=1in]{geometry}
\usepackage{amsmath,amssymb}
\usepackage{enumitem}
\usepackage{fancyhdr}
\usepackage{xcolor} 

% ===================== HEADER & FOOTER =====================
\pagestyle{fancy}
\fancyhf{}
\lhead{CSE 400: Fundamentals of Probability in Computing}
\rhead{L16 Scribe}
\cfoot{\thepage}

% ===================== CUSTOM COMMANDS =====================
\newcommand{\solution}{
    \vspace{0.5cm}
    \noindent\textbf{\textcolor{blue}{Solution:}}
    \vspace{0.2cm}
}

% ===================== TITLE =====================
\title{
    \normalsize School of Engineering and Applied Science (SEAS), Ahmedabad University \\
    \vspace{0.2cm}
    \textbf{CSE 400: Fundamentals of Probability in Computing}\\
    \Large Lecture scribe submission
}
\author{}
\date{}

\begin{document}
\maketitle

\vspace{-2cm}
\begin{center}
    \begin{tabular}{ll}
        \textbf{Name:} & Ishika Patel \\[1.5ex]
        \textbf{Enrollment No:} & AU2440088 \\[1.5ex]
        \textbf{Group:} & s1\_g1\_fin \\[1.5ex]
        \textbf{Email:} & ishika.p3@ahduni.edu.in \\[1.5ex]
        \textbf{Lecture:} & L15--L18 \\[1.5ex]
        \textbf{Date of Submission:} & March 17\textsuperscript{th}, 2026 (11:59 PM)
    \end{tabular}
\end{center}

\hrule
\vspace{0.5cm}

\section{Lecture Scribe}

\textbf{CSE400 – Fundamentals of Probability in Computing}

\textbf{Lecture L15--L18: Joint Random Variables, Joint Distributions, and Examples}

Instructor: Dhaval Patel, PhD

SEAS – Ahmedabad University

Date: March 10, 2026

\section{Outline}

The lecture covers the following topics:

\begin{itemize}
\item Why Study Joint Random Variables?
\item Discrete Case
\item Joint PMF Table Representation
\item Example: Rolling Two Dice
\item Continuous Case
\item Geometric Interpretation
\item Worked Example
\item Joint Cumulative Distribution Function
\item Definition of Joint CDF
\item Continuous Case
\item Discrete Case
\item Properties of the Joint CDF
\item Comprehensive Example
\end{itemize}

\section{Why Study Joint Random Variables?}

In science and in real life, we are often interested in two or more random variables at the same time.

\subsection{Example 1: Smart Transportation}

Two random variables:

\[
X = \text{Vehicle Speed}
\]

\[
Y = \text{Traffic Density}
\]

Question:

Can speed be high when the traffic density is also high?

Answer:

Yes / No

Conclusion:

X and Y are correlated random variables.

Real ITS Applications (probabilistic analysis):

\begin{itemize}
\item Adaptive traffic signals
\item Congestion prediction
\item Autonomous vehicle routing
\end{itemize}

\subsection{Example 2: Wireless Network Performance}

Two random variables:

\[
X = \text{Signal-to-Noise Ratio (SNR)}
\]

\[
Y = \text{Data Throughput}
\]

Reason for studying jointly:

Throughput depends on SNR.

Both vary due to:

\begin{itemize}
\item fading
\item interference
\item congestion
\end{itemize}

Question:

If SNR is high, will throughput always be high?

Answer:

Yes / No

Reasoning:

Network congestion creates conditional relationships.

Applications in networking:

\begin{itemize}
\item 5G scheduling
\item Adaptive modulation
\item Network planning
\end{itemize}

\subsection{Example 3: Weather Prediction}

Random variables:

\[
X = \text{Rainfall}
\]

\[
Y = \text{Temperature}
\]

Application:

Weather forecasting.

\subsection{Example 4: Drone Delivery}

Random variables:

\[
X = \text{Wind Speed}
\]

\[
Y = \text{Delivery Time}
\]

Application:

Logistics and UAV systems.

\subsection{Example 5: Rolling Two Dice}

Two random variables:

\[
X = \text{Outcome of first die}
\]

\[
Y = \text{Outcome of second die}
\]

Example events:

\[
X + Y = 7
\]

\[
X > Y
\]

This requires the joint PMF.

Many (infinite) examples exist where multiple random variables must be studied together.

\section{Discrete Case – Joint PMF}

The discrete case introduces the concept of the Joint Probability Mass Function (Joint PMF) for two discrete random variables.

\subsection{Joint PMF Table Representation}

The joint PMF can be represented using a table where:

\begin{itemize}
\item Rows correspond to values of one random variable.
\item Columns correspond to values of the other random variable.
\item Each cell contains the joint probability.
\end{itemize}

\section{Example: Rolling Two Dice}

Two random variables are defined:

\[
X = \text{Outcome of die 1}
\]

\[
Y = \text{Outcome of die 2}
\]

The joint PMF describes probabilities of all combinations:

\[
P(X=x,Y=y)
\]

The joint distribution can be represented in a table form.

\section{Continuous Case – Joint PDF}

For continuous random variables, the joint distribution is described using a Joint Probability Density Function (Joint PDF).

The continuous case is interpreted geometrically.

\section{Example (Continuous Case)}

An example is provided illustrating a joint probability density function and the computation of probabilities.

The example is followed by step-by-step solutions.

\section{Joint Cumulative Distribution Function (Joint CDF)}

\subsection{Definition}

The Joint CDF for two random variables $X$ and $Y$ is defined as:

\[
F_{X,Y}(x,y) = P(X \le x, Y \le y)
\]

\subsection{Joint CDF – Continuous Case}

The joint CDF can be obtained from the joint PDF by integrating over the appropriate region.

\subsection{Joint CDF – Discrete Case}

For discrete random variables, the joint CDF is obtained by summing the joint PMF over the appropriate values.

\section{Properties of the Joint CDF}

The lecture introduces the properties of the Joint CDF.

(These properties are listed across multiple slides.)

\section{Summary of Topics Covered}

The lecture covered:

\begin{itemize}
\item Motivation for studying joint random variables
\item Joint PMF for discrete random variables
\item Joint PDF for continuous random variables
\item Table representation of joint distributions
\item Example: Rolling two dice
\item Joint Cumulative Distribution Function
\item Joint CDF in continuous and discrete cases
\item Properties of the joint CDF
\end{itemize}

\bigskip
\hrule
\vspace{0.2cm}
\begin{center}
    \small \textit{End of Submission}
\end{center}

\end{document}